\chapter{Lebih lanjut tentang masalah nilai eigen} \label{SL:chapter}

%%%%%%%%%%%%%%%%%%%%%%%%%%%%%%%%%%%%%%%%%%%%%%%%%%%%%%%%%%%%%%%%%%%%%%%%%%%%%%

\section{Masalah Sturm--Liouville}
\label{slproblems:section}

%mbxINTROSUBSECTION

\sectionnotes{2 kuliah\EPref{, \S10.1 dalam \cite{EP}}\BDref{,
\S11.2 dalam \cite{BD}}}

\subsection{Masalah nilai batas}

Dalam \chapterref{FS:chapter},
kita menjumpai beberapa masalah nilai eigen yang berbeda, seperti:
\begin{equation*}
X''(x) + \lambda X(x) = 0 ,
\end{equation*}
dengan syarat batas yang berbeda
\index{syarat batas Dirichlet}%
\index{syarat batas Neumann}%
\index{syarat batas campuran}
\begin{equation*}
\begin{array}{rrl}
X(0) = 0 & ~~X(L) = 0 & ~~\text{(Dirichlet), atau} \\
X'(0) = 0 & ~~X'(L) = 0 & ~~\text{(Neumann), atau} \\
X'(0) = 0 & ~~X(L) = 0 & ~~\text{(Campuran), atau} \\
X(0) = 0 & ~~X'(L) = 0 & ~~\text{(Campuran)}, \ldots
\end{array}
\end{equation*}
Masalah-masalah batas ini muncul dalam kajian persamaan panas $u_t =
k u_{xx}$ ketika kita mencoba menyelesaikan persamaan tersebut dengan metode
pemisahan variabel dalam \sectionref{heateq:section}.
Syarat Dirichlet berkaitan dengan penerapan
suhu nol pada ujung-ujungnya, Neumann berarti mengisolasi ujung-ujungnya, dan seterusnya.
Jenis syarat titik ujung lainnya juga muncul secara alami, seperti
\emph{\myindex{syarat batas Robin}}
\begin{equation*}
hX(0) - X'(0) = 0, \qquad hX(L) + X'(L) = 0 ,
\end{equation*}
untuk suatu konstanta $h$.  Syarat-syarat ini muncul ketika ujung-ujung tersebut dicelupkan
ke dalam suatu medium.

Dalam perhitungan pemisahan variabel,
kita menjumpai suatu masalah nilai eigen dan menemukan
fungsi-fungsi eigen
$X_n(x)$.  Kemudian kita menemukan \emph{\myindex{dekomposisi fungsi eigen}} dari
suhu awal
$f(x) = u(x,0)$,
\begin{equation*}
f(x) = \sum_{n=1}^\infty c_n X_n(x) .
\end{equation*}
Setelah memiliki dekomposisi ini
dan menemukan $T_n(t)$ yang sesuai sedemikian sehingga $T_n(0) = 1$
dan $T_n(t)X_n(x)$ merupakan solusi persamaan panas,
kita menuliskan solusi masalah semula, termasuk syarat
awalnya, sebagai
\begin{equation*}
u(x,t) = \sum_{n=1}^\infty c_n T_n(t) X_n(x) .
\end{equation*}

\medskip

Untuk mengkaji masalah yang lebih umum dengan metode ini, kita harus mengkaji
masalah nilai eigen yang lebih umum.
Pertama, kita mempelajari 
persamaan linear orde dua berbentuk
\begin{equation} \label{SL:eq}
\frac{d}{dx} \left( p(x) \frac{dy}{dx} \right)
- q(x) y + \lambda r(x) y = 0 .
\end{equation}
Pada dasarnya,
sebarang persamaan linear orde dua
berbentuk $a(x) y'' + b(x) y' + c(x) y + \lambda d(x) y = 0$
dapat dituliskan sebagai \eqref{SL:eq}
setelah dikalikan dengan faktor yang tepat.

\begin{example}[Bessel]
Ubahlah persamaan berikut ke dalam bentuk \eqref{SL:eq}:
\begin{equation*}
x^2 y'' + xy' + \left(\lambda x^2 - n^2\right)y = 0 .
\end{equation*}
Kalikan kedua ruas dengan $\frac{1}{x}$ untuk memperoleh
\begin{equation*}
\begin{split}
\frac{1}{x} \left( x^2 y'' + xy' + \left(\lambda x^2 - n^2\right)y \right)
& =
x y'' + y' + \left(\lambda x - \frac{n^2}{x}\right)y 
\\
& =
\frac{d}{dx} \left( x \frac{dy}{dx} \right)
- \frac{n^2}{x} y + \lambda x y  = 0.
\end{split}
\end{equation*}
Persamaan Bessel muncul, misalnya, dalam penyelesaian
persamaan gelombang dua dimensi.
Jika Anda ingin melihat cara menyelesaikan persamaan tersebut,
Anda dapat melihat \subsectionref{bessel:subsection}.
\end{example}

Apa yang disebut 
\emph{\myindex{masalah Sturm--Liouville}}%
\footnote{Dinamai menurut matematikawan Prancis
\href{https://en.wikipedia.org/wiki/Jacques_Charles_Fran\%C3\%A7ois_Sturm}{Jacques Charles Fran\c{c}ois Sturm}
(1803--1855) dan
\href{https://en.wikipedia.org/wiki/Liouville}{Joseph Liouville}
(1809--1882).} adalah mencari
solusi taktrivial dari
\begin{equation} \label{sl:slprob}
\mybxbg{~~
\begin{aligned}
&\frac{d}{dx} \left( p(x) \frac{dy}{dx} \right)
- q(x) y + \lambda r(x) y = 0, \qquad a < x < b, \\
&\alpha_1 y(a) - \alpha_2 y'(a) = 0, \\
&\beta_1 y(b) + \beta_2 y'(b) = 0.
\end{aligned}
~~}
\end{equation}
Dengan taktrivial, sekali lagi kita memaksudkan suatu $y$ yang bukan sekadar
konstanta nol (yang selalu merupakan solusi).
Secara khusus, kita menanyakan nilai-nilai $\lambda$ mana
yang memungkinkan solusi taktrivial tersebut.
Nilai-nilai $\lambda$ yang mengizinkan solusi taktrivial
disebut \emph{nilai eigen\index{nilai eigen}},
dan solusi taktrivial yang bersesuaian
disebut
\emph{fungsi eigen\index{fungsi eigen}}.
Konstanta $\alpha_1$ dan $\alpha_2$
tidak boleh keduanya nol;
hal yang sama berlaku untuk
$\beta_1$ dan $\beta_2$.

\begin{theorem} \label{sl:slregthm}
Misalkan $p(x)$, $p'(x)$, $q(x)$, dan $r(x)$ kontinu pada $[a,b]$
dan misalkan $p(x) > 0$ serta $r(x) > 0$ untuk semua $x$ dalam $[a,b]$.
Maka masalah Sturm--Liouville \eqref{sl:slprob}
memiliki suatu barisan nilai eigen yang meningkat
\begin{equation*}
\lambda_1 < \lambda_2 < \lambda_3 < \cdots 
\end{equation*}
sedemikian sehingga
\begin{equation*}
\lim_{n \to \infty} \lambda_n = +\infty
\end{equation*}
dan sedemikian sehingga untuk setiap $\lambda_n$ terdapat (hingga suatu kelipatan konstanta)
satu fungsi eigen $y_n(x)$. 

Selain itu, jika $q(x) \geq 0$ dan $\alpha_1, \alpha_2, \beta_1, \beta_2 \geq 0$,
maka $\lambda_n \geq 0$ untuk semua $n$.
\end{theorem}

Masalah yang memenuhi hipotesis teorema tersebut
(termasuk bagian \myquote{Selain itu})
disebut
\emph{masalah Sturm--Liouville reguler\index{masalah Sturm--Liouville reguler}}%
\index{masalah Sturm--Liouville!reguler},
dan kita hanya akan membahas masalah seperti itu di sini.
Artinya, masalah reguler adalah masalah dengan
$p(x)$, $p'(x)$, $q(x)$, dan $r(x)$ kontinu, $p(x) > 0$, $r(x) > 0$,
$q(x) \geq 0$, dan $\alpha_1, \alpha_2, \beta_1, \beta_2 \geq 0$,
dengan tidak berlaku $\alpha_1=\alpha_2=0$
maupun
$\beta_1=\beta_2=0$.
Catatan: Berhati-hatilah terhadap tanda-tandanya.  Berhati-hatilah pula terhadap pertidaksamaan
untuk $r$ dan $p$.
Pertidaksamaan tersebut
harus tegas untuk semua $x$ dalam interval $[a,b]$, termasuk titik-titik ujungnya!

Ketika nol
merupakan suatu nilai eigen, demi kemudahan biasanya kita
mulai memberi label nilai-nilai eigen dari 0, bukan dari 1.
Artinya, kita memberi label nilai-nilai eigen $\lambda_0 < \lambda_1 < \lambda_2 < \cdots$.

\begin{example}
Masalah $y''+\lambda y=0$, $0 < x < L$, $y(0) = 0$, dan $y(L) = 0$
adalah masalah Sturm--Liouville reguler: $p(x) = 1$, $q(x) = 0$, $r(x) = 1$,
dan kita memiliki $p(x) = 1 > 0$ serta $r(x) = 1 > 0$.  Kita juga memiliki
$a=0$, $b=L$, $\alpha_1 = \beta_1 = 1$, $\alpha_2 = \beta_2 = 0$.
Nilai-nilai eigennya adalah $\lambda_n = \frac{n^2 \pi^2}{L^2}$ dan fungsi-fungsi eigennya
adalah $y_n(x) = \sin\bigl(\frac{n\pi}{L} x\bigr)$.  Semua nilai eigen nonnegatif sebagaimana
diprediksi oleh teorema tersebut.
\end{example}

\begin{exercise}
Carilah nilai eigen dan fungsi eigen untuk
\begin{equation*}
y'' + \lambda y = 0, \quad y'(0) = 0, \quad y'(1) = 0.
\end{equation*}
Identifikasikan
$p, q, r, \alpha_j, \beta_j$.
Dapatkah Anda menggunakan teorema di atas untuk mempermudah pencarian nilai eigen?
Petunjuk: Perhatikan syarat $-y'(0)=0$.
\end{exercise}

\begin{example}
Carilah nilai eigen dan fungsi eigen dari masalah
\begin{align*}
& y''+\lambda y = 0, \quad 0 < x < 1 , \\
& hy(0)- y'(0) = 0, \quad y'(1)  = 0, \quad h > 0.
\end{align*}

Persamaan-persamaan ini memberikan suatu masalah Sturm--Liouville reguler.

\begin{exercise}
Identifikasikan $p, q, r, \alpha_j, \beta_j$ dalam contoh di atas.
\end{exercise}

Menurut \thmref{sl:slregthm},
$\lambda \geq 0$.
Jadi solusi umumnya (tanpa syarat batas) adalah
\begin{equation*}
\begin{aligned}
 & y(x) = A \cos \bigl( \sqrt{\lambda}\, x\bigr)
+ B \sin \bigl( \sqrt{\lambda}\, x\bigr) & & \qquad \text{jika } \; \lambda > 0 , \\
& y(x) = A x + B & & \qquad \text{jika } \; \lambda = 0 .
\end{aligned}
\end{equation*}

Mari kita lihat apakah $\lambda = 0$ merupakan nilai eigen:
Kita harus memenuhi $0 = hB - A$ dan $A = 0$, sehingga $B=0$ (karena $h > 0$).
Oleh karena itu, 0 bukanlah nilai eigen (tidak ada solusi taknol, sehingga tidak ada fungsi eigen).

Sekarang kita
mencoba $\lambda > 0$.  Kita substitusikan syarat-syarat batasnya:
\begin{align*}
& 0 = h A - \sqrt{\lambda}\, B , \\
& 0 = -A \sqrt{\lambda}\, \sin \bigl(\sqrt{\lambda}\bigr)
+ B \sqrt{\lambda}\, \cos \bigl(\sqrt{\lambda}\bigr) .
\end{align*}
Jika $A=0$, maka $B=0$ dan sebaliknya; oleh karena itu, keduanya taknol.
Jadi $B = \frac{hA}{\sqrt{\lambda}}$, dan
$0 = -A \sqrt{\lambda}\, \sin \bigl( \sqrt{\lambda}\bigr) + \frac{hA}{\sqrt{\lambda}}
\sqrt{\lambda}\, \cos \bigl( \sqrt{\lambda}\bigr)$.  Karena $A \neq 0$, kita memperoleh
\begin{equation*}
0 = 
- \sqrt{\lambda}\, \sin \bigl( \sqrt{\lambda}\bigr)
+ h \cos \bigl( \sqrt{\lambda}\bigr) ,
\end{equation*}
atau
\begin{equation*}
\frac{h}{\sqrt{\lambda}} = \tan \sqrt{\lambda} .
\end{equation*}

Kita menggunakan komputer untuk mencari $\lambda_n$.  Tabel-tabel tersedia,
meskipun saat ini penggunaan komputer atau kalkulator grafik
jauh lebih praktis.
Metode termudah adalah menggambar grafik fungsi 
$\nicefrac{h}{x}$ dan $\tan x$, lalu melihat nilai $x$ tempat keduanya berpotongan.
Terdapat takhingga banyak titik potong.  Nyatakan
titik potong pertama
dengan $\sqrt{\lambda_1}$,
titik potong kedua dengan $\sqrt{\lambda_2}$, dan seterusnya.
Sebagai contoh, ketika
$h=1$, kita memperoleh $\sqrt{\lambda_1} \approx 0.86$, 
$\sqrt{\lambda_2} \approx 3.43$, \ldots.
Artinya, $\lambda_1 \approx 0.74$, $\lambda_2 \approx 11.73$, \ldots.
Grafik untuk $h=1$ diberikan dalam \figurevref{sl:tanx1overxfig}.
Fungsi eigen yang sesuai
(ambil $A = 1$ demi kemudahan, maka
$B=\nicefrac{h}{\sqrt{\lambda}}$) adalah
\begin{equation*}
y_n(x) = \cos \bigl( \sqrt{\lambda_n}\, x \bigr) + \frac{h}{\sqrt{\lambda_n}}
\sin \bigl(\sqrt{\lambda_n} \, x \bigr) .
\end{equation*}
Ketika $h=1$, kita memperoleh (secara hampiran)
\begin{equation*}
y_1(x) \approx \cos (0.86\, x ) + \frac{1}{0.86}
\sin (0.86 \, x ) , \qquad
y_2(x) \approx \cos (3.43\, x ) + \frac{1}{3.43}
\sin (3.43 \, x ) , \qquad \ldots .
\end{equation*}
\begin{myfig}
\capstart
\diffyincludegraphics{width=3in}{width=4.5in}{sl-tanx1overx}{%
Pada bidang xy, terdapat grafik dua fungsi di atas interval 0 hingga 7.
Fungsi pertama, 1 per x,
turun dari bagian atas gambar di sebelah kiri,
merosot sangat cepat menuju sumbu x, lalu terus mendekatinya secara perlahan tetapi
tidak pernah memotongnya ketika x membesar.  Fungsi kedua adalah tangen, yang
dimulai dari nol ketika x sama dengan 0, lalu melengkung cepat ke atas dan keluar dari gambar
sebelum pi per 2, sembari memotong grafik 1 per x untuk pertama kalinya.  Setelah pi
per 2,
tangen muncul kembali dari bagian bawah gambar, melengkung ke kanan lalu
ke kiri sebelum kembali keluar dari gambar sebelum 3 pi per 2.  Grafik tersebut memotong grafik
1 per x untuk kedua kalinya dalam interval itu.  Setelah 3 pi per 2, grafik tersebut muncul kembali dari
bagian bawah dan kembali bergerak naik serta memotong grafik 1 per x untuk ketiga kalinya.}
\caption{Grafik $\frac{1}{x}$ dan $\tan x$.%
\label{sl:tanx1overxfig}}
\end{myfig}
\end{example}

\subsection{Keortogonalan}

Kita telah melihat konsep keortogonalan sebelumnya.  Sebagai contoh,
kita telah menunjukkan bahwa $\sin (nx)$ saling ortogonal untuk $n$ yang berbeda pada $[0,\pi]$.
Untuk masalah Sturm--Liouville umum, kita memerlukan kerangka yang lebih umum.
Misalkan $r(x)$
merupakan suatu \emph{\myindex{fungsi bobot}} (sebarang fungsi, meskipun secara umum kita
mengasumsikannya positif) pada $[a,b]$.  Dua fungsi $f(x)$ dan $g(x)$
dikatakan
\emph{ortogonal\index{ortogonal!terhadap suatu bobot}} 
terhadap fungsi bobot
$r(x)$ jika
\begin{equation*}
\int_a^b f(x) \, g(x) \, r(x) \,dx = 0 .
\end{equation*}
Dalam kerangka ini,
kita mendefinisikan \emph{hasil kali dalam\index{hasil kali dalam fungsi}} sebagai
\begin{equation*}
\langle f , g \rangle \overset{\text{def}}{=} \int_a^b f(x) \, g(x) \, r(x)
\,dx ,
\end{equation*}
dan kemudian mengatakan $f$ dan $g$ ortogonal setiap kali $\langle f , g \rangle = 0$.
Hasil dan konsepnya kembali analog dengan 
aljabar linear berdimensi hingga.

Gagasan dari hasil kali dalam yang diberikan adalah bahwa nilai-nilai $x$ tempat $r(x)$
lebih besar memiliki bobot lebih besar.
Fungsi taktrivial (takkonstan)
$r(x)$ muncul secara alami, misalnya dari suatu perubahan variabel.
Oleh karena itu, Anda dapat
membayangkan suatu perubahan variabel sedemikian sehingga $d\xi = r(x)\, dx$.

Fungsi eigen dari suatu
masalah Sturm--Liouville reguler memenuhi sifat keortogonalan, sama seperti
fungsi eigen dalam \sectionref{bvp:section}.
Buktinya sangat serupa dengan bukti 
\thmvref{bvp:orthogonaleigen} yang analog.

\begin{theorem}
Misalkan kita memiliki suatu masalah Sturm--Liouville reguler
\begin{align*}
&\frac{d}{dx} \left( p(x) \frac{dy}{dx} \right)
- q(x) y + \lambda r(x) y = 0 , \\
&\alpha_1 y(a) - \alpha_2 y'(a) = 0 , \\
&\beta_1 y(b) + \beta_2 y'(b) = 0 .
\end{align*}
Misalkan $y_j$ dan $y_k$ merupakan dua fungsi eigen berbeda untuk dua
nilai eigen berbeda $\lambda_j$ dan $\lambda_k$.  Maka
\begin{equation*}
\int_a^b y_j(x) \, y_k(x) \, r(x) \,dx = 0,
\end{equation*}
artinya, $y_j$ dan $y_k$ ortogonal terhadap fungsi bobot
$r$.
\end{theorem}


\subsection{Alternatif Fredholm}

Teorema \emph{alternatif Fredholm} yang telah kita bicarakan sebelumnya
(\thmvref{thm:fredholmsimple})
berlaku untuk semua masalah Sturm--Liouville reguler.
Kita menyatakannya di sini demi kelengkapan.

\begin{theorem}[Alternatif Fredholm]%
\index{alternatif Fredholm!masalah Sturm--Liouville}
Misalkan kita memiliki suatu masalah Sturm--Liouville reguler.
Maka salah satu dari hal berikut berlaku:
\begin{align*}
&\frac{d}{dx} \left( p(x) \frac{dy}{dx} \right)
- q(x) y + \lambda r(x) y = 0 , \\
&\alpha_1 y(a) - \alpha_2 y'(a) = 0 , \\
&\beta_1 y(b) + \beta_2 y'(b) = 0 ,
\end{align*}
memiliki suatu solusi taknol ($\lambda$ merupakan nilai eigen), atau
\begin{align*}
&\frac{d}{dx} \left( p(x) \frac{dy}{dx} \right)
- q(x) y + \lambda r(x) y = f(x) , \\
&\alpha_1 y(a) - \alpha_2 y'(a) = 0 , \\
&\beta_1 y(b) + \beta_2 y'(b) = 0 ,
\end{align*}
memiliki solusi tunggal untuk setiap $f(x)$ yang kontinu pada $[a,b]$.
\end{theorem}

Teorema ini digunakan dengan cara yang hampir sama seperti sebelumnya dalam
\sectionref{sec:scs}.  Teorema ini digunakan ketika menyelesaikan masalah nilai batas
nonhomogen yang lebih umum.  Teorema tersebut tidak membantu kita menyelesaikan masalahnya.
Teorema itu memberi tahu kita kapan solusi tunggal ada, sehingga
kita mengetahui kapan perlu meluangkan waktu untuk mencarinya.  Untuk menyelesaikan masalah tersebut,
kita mendekomposisi $f(x)$ dan $y(x)$ dalam fungsi-fungsi eigen dari masalah
homogen,
kemudian menentukan koefisien deret untuk $y(x)$.

\subsection{Deret fungsi eigen}

Yang ingin kita lakukan dengan fungsi eigen setelah memperolehnya adalah
menghitung \emph{\myindex{dekomposisi fungsi eigen}} dari suatu fungsi sebarang
$f(x)$.  Artinya, kita ingin menuliskan
\begin{equation} \label{sl:fdecomp}
f(x) = \sum_{n=1}^\infty c_n y_n(x) ,
\end{equation}
dengan $y_n(x)$ merupakan fungsi-fungsi eigen.
Kita ingin mengetahui apakah kita dapat merepresentasikan
sebarang fungsi $f(x)$ dengan cara ini,
dan jika demikian, kita ingin menghitung $c_n$ (dan tentu saja kita ingin mengetahui apakah
jumlahnya konvergen).  Baiklah\@, bayangkan
kita dapat menuliskan $f(x)$ sebagai \eqref{sl:fdecomp}.  Kita akan mengasumsikan kekonvergenan dan
kemampuan untuk mengintegralkan deret suku demi suku.
Karena keortogonalan, kita mempunyai
\begin{equation*}
\begin{split}
\langle f , y_m \rangle
 = \int_a^b f(x) \, y_m (x) \, r(x) \, dx
& = \int_a^b \left( \sum_{n=1}^\infty c_n y_n(x) \right) \, y_m (x) \, r(x)
\, dx\\
&= \sum_{n=1}^\infty c_n \int_a^b y_n(x) \, y_m (x) \, r(x) \, dx\\
&= c_m \int_a^b y_m(x) \, y_m (x) \, r(x) \, dx = c_m \langle y_m , y_m \rangle
.
\end{split}
\end{equation*}
Oleh karena itu,
\begin{equation} \label{sl:cm}
\mybxbg{~~
c_m = \frac{\langle f , y_m \rangle}{\langle y_m , y_m \rangle}
=
\frac{\int_a^b f(x) \, y_m (x)\, r(x) \, dx}%
{\int_a^b {\bigl(y_m(x)\bigr)}^2 \, r(x) \,dx} .
~~}
\end{equation}

Perhatikan bahwa $y_m$ diketahui hingga suatu kelipatan konstanta, sehingga kita dapat memilih
kelipatan skalar suatu fungsi eigen sedemikian sehingga
$\langle y_m , y_m \rangle = 1$ (jika kita memiliki suatu fungsi eigen sebarang
$\tilde{y}_m$, bagilah fungsi tersebut
dengan $\sqrt{\langle \tilde{y}_m , \tilde{y}_m \rangle}$).
Ketika
$\langle y_m , y_m \rangle = 1$,
kita memperoleh bentuk yang
lebih sederhana $c_m = \langle f, y_m \rangle$.
% seperti yang kita lakukan untuk deret Fourier.
Teorema berikut berlaku
secara lebih umum, tetapi pernyataan yang diberikan memadai untuk keperluan kita.

\begin{theorem}
Misalkan $f$ merupakan fungsi kontinu mulus sepotong-sepotong pada $[a,b]$.  Jika $y_1,
y_2, \ldots$ merupakan fungsi-fungsi eigen dari suatu masalah Sturm--Liouville reguler,
masing-masing satu untuk setiap nilai eigen,
maka terdapat konstanta-konstanta real $c_1, c_2, \ldots$ yang diberikan oleh \eqref{sl:cm}
sedemikian sehingga
\eqref{sl:fdecomp} konvergen dan berlaku untuk $a < x < b$.
\end{theorem}

\begin{example}
Perhatikan
\begin{align*}
& y'' + \lambda y = 0, \quad 0 < x < \nicefrac{\pi}{2} , \\
& y(0) =0, \quad y'(\nicefrac{\pi}{2}) = 0 .
\end{align*}
Masalah di atas adalah suatu masalah Sturm--Liouville reguler, dan
\thmvref{sl:slregthm}
menyatakan bahwa jika $\lambda$ merupakan nilai eigen, maka
$\lambda \geq 0$.

Misalkan $\lambda = 0$.  Solusi umumnya adalah $y(x) = Ax + B$.
Kita substitusikan
syarat-syarat awal untuk memperoleh $0=y(0) = B$, dan $0 = y'(\nicefrac{\pi}{2}) = A$.
Oleh karena itu, $\lambda = 0$ bukanlah nilai eigen.

Jadi, mari kita perhatikan $\lambda > 0$, dengan
solusi umum
\begin{equation*}
y(x) = A \cos \bigl( \sqrt{\lambda} \, x \bigr)
+ B \sin \bigl( \sqrt{\lambda} \, x \bigr) .
\end{equation*}
Dengan menyubstitusikan syarat-syarat batas, kita memperoleh
$0 = y(0) = A$ dan $0 = y'(\nicefrac{\pi}{2})
= \sqrt{\lambda} \, B \cos \bigl(\sqrt{\lambda} \, \frac{\pi}{2}\bigr)$.
Karena $A$ nol,
$B$ tidak boleh nol.  Oleh karena itu, $\cos \bigl( \sqrt{\lambda} \,
\frac{\pi}{2}\bigr) = 0$.
Ini berarti bahwa
$\sqrt{\lambda} \,\frac{\pi}{2}$ merupakan kelipatan bilangan bulat ganjil dari
$\nicefrac{\pi}{2}$,
yaitu $(2n-1)\frac{\pi}{2} = \sqrt{\lambda_n} \,\frac{\pi}{2}$.
Dengan menyelesaikannya untuk $\lambda_n$, kita memperoleh
\begin{equation*}
\lambda_n = {(2n-1)}^2 .
\end{equation*}
Kita dapat mengambil $B = 1$.  Fungsi-fungsi eigen kita adalah
\begin{equation*}
y_n(x) = \sin \bigl( (2n-1)x \bigr) .
\end{equation*}
Sedikit kalkulus menunjukkan
\begin{equation*}
\int_0^{\frac{\pi}{2}} {\Bigl( \sin \bigl( (2n-1)x \bigr) \Bigr)}^2 \, dx
= \frac{\pi}{4} .
\end{equation*}

Jadi, sebarang fungsi mulus sepotong-sepotong $f(x)$ pada $[0,\nicefrac{\pi}{2}]$ dapat dituliskan sebagai
\begin{equation*}
f(x) = \sum_{n=1}^\infty c_n \sin \bigl( (2n-1)x \bigr) ,
\end{equation*}
dengan
\begin{equation*}
c_n = \frac{\langle f , y_n \rangle}{\langle y_n , y_n \rangle}
= \frac{\int_0^{\frac{\pi}{2}} f(x) \, \sin \bigl( (2n-1)x \bigr) \, dx
}{\int_0^{\frac{\pi}{2}} {\Bigl(\sin \bigl((2n-1)x\bigr)\Bigr)}^2 \, dx}
= \frac{4}{\pi} \int_0^{\frac{\pi}{2}} f(x) \,\sin \bigl( (2n-1)x \bigr) \, dx .
\end{equation*}

Perhatikan bahwa deret tersebut konvergen ke perluasan ganjil yang periodik-$2\pi$
dari $f(x)$.  Dengan deret sinus reguler, kita akan mengharapkan
fungsi berperiode $2 \, \frac{\pi}{2} = \pi$.

\begin{exercise}[menantang]
Dalam contoh di atas, fungsi tersebut didefinisikan pada $0 < x < \nicefrac{\pi}{2}$,
namun deret terhadap fungsi-fungsi eigen
$\sin \bigl( (2n-1)x \bigr)$ konvergen ke perluasan ganjil yang periodik-$2\pi$ dari $f(x)$.
Tentukan bagaimana perluasan tersebut didefinisikan untuk $\nicefrac{\pi}{2} < x < \pi$.
\end{exercise}

Mari kita hitung suatu contoh.
Perhatikan $f(x) = x$ untuk $0 < x <  \nicefrac{\pi}{2}$.  Setelah beberapa
perhitungan kalkulus, kita memperoleh
\begin{equation*}
c_n = 
\frac{4}{\pi} \int_0^{\frac{\pi}{2}} f(x) \,\sin \bigl( (2n-1)x \bigr) \, dx 
=
\frac{4{(-1)}^{n+1}}{\pi {(2n-1)}^2} ,
\end{equation*}
sehingga untuk $x$ dalam $[0,\nicefrac{\pi}{2}]$,
\begin{equation*}
f(x) = \sum_{n=1}^\infty \frac{4{(-1)}^{n+1}}{\pi {(2n-1)}^2}
\sin \bigl( (2n-1)x \bigr) .
\end{equation*}
Deret ini berbeda dari deret sinus reguler periodik-$\pi$ yang dapat
dihitung sebagai
\begin{equation*}
f(x) = \sum_{n=1}^\infty \frac{{(-1)}^{n+1}}{n}  \sin ( 2nx ) .
\end{equation*}
Kedua deret konvergen ke $f(x)$ untuk $0 < x < \nicefrac{\pi}{2}$, tetapi
fungsi-fungsi eigen yang terlibat berasal dari masalah nilai eigen yang berbeda.
\end{example}

\subsection{Latihan}

\begin{exercise}
Carilah nilai eigen dan fungsi eigen dari
\begin{equation*}
y''+\lambda y = 0, \quad y(0)- y'(0) = 0, \quad y(1) = 0 .
\end{equation*}
\end{exercise}

\begin{exercise}
Kembangkan fungsi $f(x) = x$ pada $0 \leq x \leq 1$ dengan menggunakan fungsi-fungsi eigen
dari sistem
\begin{equation*}
y'' + \lambda y = 0, \quad y'(0) = 0, \quad y(1) = 0 .
\end{equation*}
\end{exercise}

\begin{exercise}
Misalkan Anda memiliki suatu masalah Sturm--Liouville pada interval
$[0,1]$ dan memperoleh
$y_n(x) = \sin (\gamma n x)$, dengan $\gamma > 0$ suatu konstanta.
Dekomposisikan $f(x) = x$, $0 < x < 1$, dalam fungsi-fungsi eigen tersebut.
\end{exercise}

\begin{exercise}
Carilah nilai eigen dan fungsi eigen dari
\begin{equation*}
y^{(4)}+\lambda y = 0, \quad y(0) = 0, \quad y'(0) = 0, \quad y(1) = 0, \quad
y'(1) = 0 .
\end{equation*}
Masalah ini bukan masalah Sturm--Liouville, tetapi gagasannya sama.
\end{exercise}

\begin{exercise}[lebih menantang]
Carilah nilai eigen dan fungsi eigen untuk
\begin{equation*}
\frac{d}{dx} (e^x y') + \lambda e^x y = 0, \quad y(0) = 0, \quad y(1) = 0 .
\end{equation*}
%Atau setidaknya susunlah integral-integral yang perlu diselesaikan.
Petunjuk: Pertama, tuliskan sistem tersebut sebagai sistem berkoefisien konstan untuk mencari
solusi umum.  Perhatikan bahwa \thmvref{sl:slregthm} menjamin $\lambda \geq 0$.
\end{exercise}

\setcounter{exercise}{100}

\begin{exercise}
Carilah nilai eigen dan fungsi eigen dari
\begin{equation*}
y'' + \lambda y = 0, \quad y(-1) = 0, \quad y(1) = 0 .
\end{equation*}
\end{exercise}
\exsol{%
$\lambda_n = \frac{n^2\pi^2}{4}$, $n=1,2,3,\ldots$.
Untuk $n$ ganjil,
$y_n = \cos\left(\frac{n\pi}{2} x \right)$,
untuk $n$ genap,
$y_n = \sin\left(\frac{n\pi}{2} x \right)$.
}

\begin{exercise}
Ubahlah masalah-masalah berikut ke dalam bentuk standar untuk masalah
Sturm--Liouville, yaitu carilah $p(x)$, $q(x)$, $r(x)$,
$\alpha_1$,
$\alpha_2$,
$\beta_1$, dan
$\beta_2$,
serta tentukan apakah masalah-masalah tersebut reguler atau tidak.
\begin{tasks}
\task $x y'' + \lambda y = 0$
\enspace untuk $0 < x < 1$,
\enspace $y(0) = 0$,
\enspace $y(1) = 0$.
\task
$(1+x^2) y'' + 2xy' + (\lambda-x^2) y = 0$
\enspace untuk $-1 < x < 1$,
\enspace $y(-1) = 0$,
\enspace $y(1)+y'(1) = 0$.\footnote{%
Dalam versi sebelumnya dari buku ini, suatu salah ketik membuat persamaannya
menjadi $(1+x^2) y'' - 2xy' + (\lambda-x^2) y = 0$, yang menghasilkan sesuatu
yang lebih sulit daripada yang dimaksudkan.  Cobalah persamaan ini sebagai tantangan tambahan.}
\end{tasks}
\end{exercise}
\exsol{%
a)~$p(x) = 1$, $q(x) = 0$, $r(x) = \frac{1}{x}$, $\alpha_1 = 1$, $\alpha_2 =
0$, $\beta_1 = 1$, $\beta_2 = 0$.  Masalah tersebut tidak reguler.
b)~$p(x) = 1+x^2$, $q(x) = x^2$, $r(x) = 1$, $\alpha_1 = 1$, $\alpha_2 =
0$, $\beta_1 = 1$, $\beta_2 = 1$.  Masalah tersebut reguler.
}

%%%%%%%%%%%%%%%%%%%%%%%%%%%%%%%%%%%%%%%%%%%%%%%%%%%%%%%%%%%%%%%%%%%%%%%%%%%%%%

\sectionnewpage
\section{Masalah nilai eigen orde tinggi}
\label{sec:appeig}

\sectionnotes{1 kuliah\EPref{, \S10.2 dalam \cite{EP}}\BDref{,
latihan dalam \S11.2 dari \cite{BD}}}

Deret fungsi eigen dapat muncul bahkan dari persamaan orde tinggi.
Perhatikan suatu balok elastis (misalnya terbuat dari baja).  Kita akan mengkaji
\emph{\myindex{getaran transversal}}
(getaran dengan balok tersebut \myquote{melentur})
dari balok tersebut.
Misalkan balok terletak sepanjang
sumbu-$x$ dan misalkan $y(x,t)$ mengukur simpangan titik $x$
pada balok pada waktu $t$.  Lihat \figurevref{appeig:transbeamfig}.

\begin{myfig}
\capstart
\inputpdft{trans-beam}{%
Diagram sebuah balok kaku yang diletakkan sepanjang sumbu x pada bidang xy, dengan salah satu
ujungnya di titik asal dan ujung kedua di suatu tempat pada bagian positif
sumbu x.  Balok tersebut melentur ke atas ke arah y, yang menunjukkan bahwa
grafik y memberikan bentuk balok yang melentur.}
\caption{Getaran transversal suatu balok.\label{appeig:transbeamfig}}
\end{myfig}

Persamaan yang mengatur keadaan ini adalah
\begin{equation*}
a^4 \frac{\partial^4 y}{\partial x^4} + \frac{\partial^2 y}{\partial t^2} = 0,
\end{equation*}
untuk suatu konstanta $a > 0$.
Kita tidak perlu mengkhawatirkan fisikanya\footnote{Jika
Anda tertarik, $a^4 = \frac{EI}{\rho}$, dengan $E$ merupakan modulus
elastisitas, $I$ merupakan momen luas kedua dari penampang,
dan $\rho$ merupakan kerapatan linear.}.

Misalkan balok tersebut memiliki panjang 1 dan ditumpu sederhana (bersendi) pada kedua ujungnya.
Balok tersebut disimpangkan menurut suatu fungsi $f(x)$ pada waktu $t=0$, kemudian
dilepaskan (kecepatan awalnya 0).  Maka $y$ memenuhi:
\begin{equation} \label{appeig:beameq}
\begin{aligned}
& a^4 y_{xxxx} + y_{tt} = 0 \qquad (0 < x < 1, \enspace t > 0), \\
& y(0,t) = y_{xx}(0,t) = 0 , \\
& y(1,t) = y_{xx}(1,t) = 0 , \\
& y(x,0) = f(x), \qquad y_{t}(x,0) = 0 .
\end{aligned}
\end{equation}

Sekali lagi, kita mencoba $y(x,t) = X(x)T(t)$ dan menyubstitusikannya untuk memperoleh
$a^4 X^{(4)}T + XT'' = 0$ atau 
\begin{equation*}
\frac{X^{(4)}}{X} = \frac{- T''}{a^4T} = \lambda .
\end{equation*}
Persamaan-persamaannya adalah
\begin{equation*}
T'' + \lambda a^4 T = 0
\qquad \text{dan} \qquad
X^{(4)} - \lambda X = 0 .
\end{equation*}
Syarat-syarat batas
$y(0,t) = y_{xx}(0,t) = 0$ dan $y(1,t) = y_{xx}(1,t) = 0$ mengimplikasikan
\begin{equation*}
X(0) =  X''(0) = 0
\qquad \text{dan} \qquad
X(1) =  X''(1) = 0 .
\end{equation*}
Syarat homogen awal $y_t(x,0) = 0$ mengimplikasikan
\begin{equation*}
T'(0) = 0 .
\end{equation*}
Seperti biasa, kita menunda syarat nonhomogen $y(x,0) = f(x)$ hingga nanti.

Dengan memperhatikan persamaan untuk $T$, yaitu
$T'' + \lambda a^4 T = 0$, intuisi fisik membawa kita
pada fakta bahwa jika $\lambda$ merupakan nilai eigen, maka $\lambda > 0$:
Kita mengharapkan getaran, bukan pertumbuhan atau
peluruhan eksponensial dalam arah $t$ (misalnya, tidak ada gesekan dalam model kita).
Jadi, tidak ada nilai eigen negatif.
Demikian pula, $\lambda = 0$ bukanlah nilai eigen.

\begin{exercise}
Benarkan $\lambda > 0$ hanya dari persamaan untuk $X$ dan syarat-syarat
batas.
\end{exercise}

Misalkan $\omega = \sqrt[4]{\lambda}$, yaitu $\omega^4 = \lambda$,
agar kita tidak perlu terus-menerus menuliskan akar pangkat empat.
Perhatikan bahwa $\omega > 0$.
Solusi umum untuk
$X^{(4)} - \omega^4 X = 0$ adalah
\begin{equation*}
X(x) = A e^{\omega x} + B e^{-\omega x} + C \sin (\omega x) +
D \cos (\omega x) .
\end{equation*}
Sekarang $0 = X(0) = A+B+D$, $0 = X''(0) = \omega^2 (A + B - D)$.
Dengan menyelesaikannya, kita memperoleh
$D = 0$ dan $B = -A$.
Jadi
\begin{equation*}
X(x) = A e^{\omega x} - A e^{-\omega x} + C \sin (\omega x) .
\end{equation*}
Selain itu, $0 = X(1) = A (e^{\omega} - e^{-\omega}) + C \sin \omega$, dan
$0 = X''(1) = A \omega^2 (e^{\omega} - e^{-\omega}) - C \omega^2 \sin \omega$.
Ini berarti $C \sin \omega  = 0$ dan 
$A (e^{\omega} - e^{-\omega}) = 2 A \sinh \omega = 0$.  Jika $\omega > 0$, maka
$\sinh \omega \neq 0$, sehingga $A = 0$.  Jadi $C \neq 0$; jika tidak,
$\lambda$ bukan
nilai eigen.  Selain itu, $\omega$ harus merupakan kelipatan bilangan bulat dari
$\pi$.  Oleh karena itu, $\omega = n \pi$ dan $n \geq 1$ (karena $\omega > 0$).  Kita dapat mengambil
$C=1$.  Jadi nilai-nilai eigennya adalah $\lambda_n = n^4 \pi^4$, dan fungsi-fungsi eigen yang bersesuaian
adalah $\sin (n \pi x)$.

Sekarang 
$T'' + n^4 \pi^4 a^4 T = 0$.  Solusi umumnya adalah $T(t) =
A \sin (n^2 \pi^2 a^2 t) + B \cos (n^2 \pi^2 a^2 t)$.  Namun $T'(0) = 0$, sehingga
$A=0$.  Kita mengambil $B=1$ agar $T(0) = 1$ demi kemudahan.
Jadi solusi kita adalah $T_n(t) = \cos (n^2 \pi^2 a^2 t)$.

Fungsi-fungsi eigennya hanyalah sinus, sehingga kita mendekomposisi fungsi $f(x)$
menggunakan deret sinus.  Artinya,
kita mencari bilangan-bilangan $b_n$ sedemikian sehingga untuk
$0 < x < 1$,
\begin{equation*}
f(x) = \sum_{n=1}^\infty b_n \sin (n \pi x) .
\end{equation*}
Kemudian solusi untuk \eqref{appeig:beameq} adalah
\begin{equation*}
y(x,t) = \sum_{n=1}^\infty b_n
X_n(x) T_n(t)
= \sum_{n=1}^\infty b_n
\sin (n \pi x)  \cos ( n^2 \pi^2 a^2 t ) .
\end{equation*}
Pokoknya adalah bahwa $X_nT_n$ merupakan suatu solusi yang memenuhi semua syarat
homogen (semua syarat kecuali posisi awal).  Dan karena
$T_n(0) = 1$,
\begin{equation*}
y(x,0) = \sum_{n=1}^\infty b_n X_n(x) T_n(0) = 
\sum_{n=1}^\infty b_n X_n(x) =
\sum_{n=1}^\infty b_n
\sin (n \pi x) = f(x) .
\avoidbreak
\end{equation*}
Jadi $y(x,t)$ menyelesaikan \eqref{appeig:beameq}.

Frekuensi alami (sudut) sistem tersebut adalah $n^2 \pi^2 a^2$.
Semua frekuensi ini merupakan kelipatan bilangan bulat dari frekuensi dasar
$\pi^2 a^2$, sehingga kita memperoleh nada musik yang merdu.  Frekuensi yang tepat
beserta amplitudonya
adalah apa yang disebut para musisi sebagai \emph{\myindex{timbre}} nada tersebut (di luar
musik, hal itu disebut spektrum).

Timbre suatu balok yang ditumpu sederhana
berbeda dari timbre dawai yang bergetar, karena pada dawai kita memperoleh \myquote{lebih banyak}
frekuensi rendah sebab kita memperoleh semua kelipatan bilangan bulat,
$1,2,3,4,5,\ldots$.  Untuk balok yang ditumpu sederhana, kita memperoleh
hanya kelipatan kuadrat $1,4,9,16,25,\ldots$, sehingga
kita memperoleh sesuatu yang jauh lebih dekat dengan bunyi murni.
Oleh karena itu, bunyi xilofon atau vibrafon sangat berbeda dari
bunyi gitar atau piano.

\begin{example}
Perhatikan $f(x) = \frac{x(x-1)}{10}$.  
Pada $0 < x < 1$ (sekarang Anda sudah mengetahui cara melakukannya),
\begin{equation*}
f(x) = \sum_{\substack{n=1\\n \text{~ganjil}}}^\infty \frac{-4}{5\pi^3 n^3}
\sin (n \pi x) .
\end{equation*}
Oleh karena itu, solusi untuk \eqref{appeig:beameq} dengan posisi awal
$f(x)$ yang diberikan adalah
\begin{equation*}
y(x,t) = \sum_{\substack{n=1\\n \text{~ganjil}}}^\infty \frac{-4}{5\pi^3 n^3}
\sin (n \pi x) \cos ( n^2 \pi^2 a^2 t ) .
\end{equation*}
\end{example}

Terdapat syarat batas lain selain ujung-ujung bersendi.  Ada
tiga kemungkinan dasar: bersendi, bebas, atau tetap.  Perhatikan
ujung di $x=0$.  Untuk ujung lainnya, gagasannya sama.
Jika ujung tersebut \emph{bersendi}\index{ujung bersendi pada balok}, maka
\begin{equation*}
y(0,t) = y_{xx}(0,t) = 0 .
\end{equation*}
Jika ujung tersebut \emph{bebas}\index{ujung bebas pada balok}, yaitu ujung tersebut sekadar
melayang di udara, maka
\begin{equation*}
y_{xx}(0,t) = y_{xxx}(0,t) = 0 .
\end{equation*}
Dan akhirnya, jika ujung tersebut
\emph{dijepit}\index{ujung dijepit pada balok}
atau  
\emph{tetap}\index{ujung tetap pada balok}, misalnya ujung tersebut dilas pada suatu
dinding,
maka
\begin{equation*}
y(0,t) = y_{x}(0,t) = 0 .
\end{equation*}


\subsection{Latihan}

\begin{exercise}
Misalkan Anda memiliki suatu balok dengan panjang 5 dan kedua ujungnya bebas.  Misalkan $y$ adalah
simpangan transversal balok pada posisi $x$ di sepanjang balok ($0 < x < 5$).
Anda mengetahui bahwa
konstanta-konstantanya sedemikian sehingga fungsi ini memenuhi persamaan $y_{tt} + 4 y_{xxxx} =
0$.  Misalkan Anda mengetahui bahwa bentuk awal balok adalah grafik
$x(5-x)$, dan kecepatan awalnya secara seragam sama dengan 2 (sama untuk setiap $x$)
dalam arah $y$ positif.  Susunlah persamaan beserta
syarat batas dan syarat awalnya.  Cukup susun, jangan selesaikan.
\end{exercise}

\begin{exercise}
Misalkan Anda memiliki suatu balok dengan panjang 5, dengan satu ujung bebas dan satu ujung tetap
(ujung tetap berada di $x=5$).
Misalkan $u$ adalah
simpangan longitudinal balok pada posisi $x$ di sepanjang balok ($0 < x < 5$).
Anda mengetahui bahwa
konstanta-konstantanya sedemikian sehingga fungsi ini memenuhi persamaan $u_{tt} = 4 u_{xx}$.
Misalkan Anda mengetahui bahwa simpangan awal balok
adalah $\frac{x-5}{50}$, dan kecepatan awalnya $\frac{-(x-5)}{100}$
dalam arah $u$ positif.  Susunlah persamaan beserta
syarat batas dan syarat awalnya.  Cukup susun, jangan selesaikan.
\end{exercise}

\begin{exercise}
Misalkan panjang balok adalah $L$ satuan, sementara segala hal lainnya tetap sama
seperti dalam \eqref{appeig:beameq}.  Apakah persamaan dan
solusi deretnya?
\end{exercise}

\begin{exercise}
Misalkan Anda memiliki 
\begin{equation*}
\begin{aligned}
& a^4 y_{xxxx} + y_{tt} = 0 \quad (0 < x < 1, t > 0) , \\
& y(0,t) = y_{xx}(0,t) = 0,\\
& y(1,t) = y_{xx}(1,t) = 0 ,\\
& y(x,0) = f(x), \quad y_{t}(x,0) = g(x) .
\end{aligned}
\end{equation*}
Artinya, Anda juga memiliki kecepatan awal.  Carilah suatu solusi deret.  Petunjuk:
Gunakan gagasan yang sama seperti yang kita gunakan untuk persamaan gelombang.
\end{exercise}

\setcounter{exercise}{100}

\begin{exercise}
Misalkan Anda memiliki suatu balok dengan panjang 1 dan kedua ujungnya bersendi.  Misalkan $y$ adalah
simpangan transversal balok pada posisi $x$ di sepanjang balok ($0 < x < 1$).
Anda mengetahui bahwa
konstanta-konstantanya sedemikian sehingga fungsi ini memenuhi persamaan $y_{tt} + 4 y_{xxxx} =
0$.  Misalkan Anda mengetahui bahwa bentuk awal balok adalah grafik
$\sin (\pi x)$, dan kecepatan awalnya 0.  Tentukan $y$.
\end{exercise}
\exsol{%
$y(x,t) = \sin(\pi x) \cos (2 \pi^2 t)$
}

\begin{exercise}
Misalkan Anda memiliki suatu balok dengan panjang 10 dan kedua ujungnya tetap.  Misalkan $y$ adalah
simpangan transversal balok pada posisi $x$ di sepanjang balok ($0 < x < 10$).
Anda mengetahui bahwa
konstanta-konstantanya sedemikian sehingga fungsi ini memenuhi persamaan $y_{tt} + 9 y_{xxxx} =
0$.  Misalkan Anda mengetahui bahwa bentuk awal balok adalah grafik
$\sin^2(\pi x)$, dan kecepatan awalnya sama dengan $x(10-x)$.
Susunlah persamaan beserta
syarat batas dan syarat awalnya.  Cukup susun, jangan selesaikan.
\end{exercise}
\exsol{%
$9 y_{xxxx} + y_{tt} = 0 \quad (0 < x < 10, t > 0)$, \quad
$y(0,t) = y_{x}(0,t) = 0$, \quad
$y(10,t) = y_{x}(10,t) = 0$, \quad
$y(x,0) = \sin^2(\pi x), \quad y_{t}(x,0) = x(10-x)$.
}

%%%%%%%%%%%%%%%%%%%%%%%%%%%%%%%%%%%%%%%%%%%%%%%%%%%%%%%%%%%%%%%%%%%%%%%%%%%%%%

\sectionnewpage
\section{Solusi periodik tunak}
\label{sps:section}

%mbxINTROSUBSECTION

\sectionnotes{1--2 kuliah\EPref{, \S10.3 dalam \cite{EP}}\BDref{,
tidak terdapat dalam \cite{BD}}}

\subsection{Dawai bergetar paksa}

Perhatikan suatu dawai gitar dengan panjang $L$.  Kita telah mengkaji 
keadaan ini dalam \sectionref{we:section}.
Misalkan $x$ adalah posisi pada dawai, $t$ adalah waktu, dan $y$ adalah simpangan dawai.  Lihat
\figurevref{sps:vibstrfig}.

\begin{myfig}
\capstart
\inputpdft{sps-vibstr}{%
%perhatikan bahwa diagram ini muncul lebih dari sekali dalam buku
Diagram sebuah dawai gitar (sebuah gitar berarsir berada di latar belakang),
penghubung gitar berada di x sama dengan 0 pada sumbu x, dan nut gitar berada di
x sama dengan L.  Dawai tersebut ditetapkan pada sumbu x di kedua titik ini.
Dawai tampak dipetik ke atas di bagian tengah sehingga membentuk suatu
segitiga lebar.  Sumbu vertikal diberi label y, dan sebuah panah menunjukkan bahwa
y memberikan simpangan dawai pada setiap x tertentu, yaitu
grafik y merupakan bentuk dawai tersebut.}
\caption{Dawai bergetar.\label{sps:vibstrfig}}
\end{myfig}

Masalah tersebut diatur oleh persamaan gelombang
\begin{equation} \label{sps:freevib}
\begin{array}{ll}
y_{tt} = a^2 y_{xx} , & \\
\noalign{\smallskip}
y(0,t) = 0 , & y(L,t) = 0 , \\
y(x,0) = f(x) , & y_t(x,0) = g(x) .
\end{array}
\end{equation}
Kita menemukan bahwa solusinya berbentuk
\begin{equation*}
y = 
\sum_{n=1}^\infty \left( A_n \cos \left( \frac{n\pi a}{L} t \right) +
B_n \sin \left( \frac{n\pi a}{L} t \right) \right)
\sin \left( \frac{n\pi}{L} x \right) ,
\end{equation*}
dengan $A_n$ dan $B_n$ ditentukan oleh syarat-syarat awal.  Frekuensi alami
dari sistem tersebut adalah frekuensi (sudut) $\frac{n \pi a}{L}$
untuk bilangan bulat $n \geq 1$.

Akan tetapi, ini adalah getaran bebas.  Bagaimana jika ada gaya luar yang bekerja pada
dawai?  Misalkan, sebagai contoh, terdapat getaran udara (kebisingan), misalnya dari dawai
kedua.  Atau mungkin dari mesin jet.  Demi kesederhanaan, asumsikan bunyi murni yang baik
dan asumsikan gaya tersebut seragam pada setiap posisi di dawai:
Gaya luar sebagai fungsi waktu $t$ diberikan oleh
$F_0 \cos (\omega t)$ sebagai gaya per satuan massa.
Kita akan menyebutnya fungsi pemaksa.  Maka persamaan
gelombang kita menjadi (ingat bahwa gaya sama dengan massa dikalikan percepatan)
\begin{equation} \label{sps:forcedeq}
y_{tt} = a^2 y_{xx} + F_0 \cos ( \omega t) ,
\end{equation}
tentu saja dengan syarat batas yang sama.

Misalkan kita ingin mencari solusi di sini yang memenuhi persamaan
\eqref{sps:forcedeq} dan
\begin{equation} \label{sps:forcedinitcond}
y(0,t) = 0, \qquad y(L,t) = 0, \qquad
y(x,0) = 0, \qquad y_t(x,0) = 0.
\end{equation}
Artinya, dawai tersebut mula-mula diam.  Pertama, kita mencari suatu solusi khusus
$y_p$ dari \eqref{sps:forcedeq} yang hanya memenuhi
syarat-syarat
$y(0,t) = y(L,t) = 0$.  Kita mendefinisikan fungsi $f$ dan $g$ sebagai
\begin{equation*}
f(x) = -y_p(x,0), \qquad g(x) = -\frac{\partial y_p}{\partial t} (x,0) .
\end{equation*}
Kemudian kita mencari suatu solusi $y_c$ dari \eqref{sps:freevib}.  Jika kita menjumlahkan kedua
solusi tersebut, kita memperoleh bahwa $y = y_c + y_p$ menyelesaikan \eqref{sps:forcedeq} dengan
syarat-syarat awal tersebut.

\begin{exercise}
Periksa bahwa $y = y_c + y_p$ menyelesaikan \eqref{sps:forcedeq} dan
syarat-syarat samping \eqref{sps:forcedinitcond}.
\end{exercise}

Jadi, persoalan utama di sini adalah mencari solusi khusus $y_p$.
Kita memperhatikan persamaannya dan membuat tebakan yang beralasan
\begin{equation*}
y_p(x,t) = X(x) \cos (\omega t) .
\end{equation*}
Kita menyubstitusikannya untuk memperoleh
\begin{equation*}
-\omega^2 X \cos ( \omega t) = a^2 X'' \cos ( \omega t) +
F_0 \cos ( \omega t ) ,
\end{equation*}
atau
$-\omega^2 X = a^2 X'' + F_0$ setelah mencoret kosinus.
Kita mengetahui cara mencari solusi umum persamaan ini (persamaan ini adalah
persamaan nonhomogen berkoefisien konstan).
Solusi umumnya adalah
\begin{equation*}
X(x) = A \cos \left( \frac{\omega}{a} x \right)
+ B \sin \left( \frac{\omega}{a} x \right) -
\frac{F_0}{\omega^2} .
\end{equation*}
Syarat-syarat titik ujung mengimplikasikan $X(0) = X(L) = 0$.  Jadi
\begin{equation*}
0 = X(0) = A - \frac{F_0}{\omega^2} ,
\end{equation*}
atau $A = \frac{F_0}{\omega^2}$, dan juga
\begin{equation*}
0 = X(L)
= \frac{F_0}{\omega^2} \cos \left( \frac{\omega L}{a} \right)
+ B \sin \left( \frac{\omega L}{a} \right) -
\frac{F_0}{\omega^2} .
\end{equation*}
Dengan mengasumsikan bahwa $\sin ( \frac{\omega L}{a} )$ tidak nol, kita dapat menyelesaikannya untuk $B$ dan
memperoleh
\begin{equation} \label{natfreq:Beq}
B = 
\frac{-F_0 \left( \cos \left( \frac{\omega L}{a} \right) - 1 \right)}%
{\omega^2 \sin \left( \frac{\omega L}{a} \right)}.
\end{equation}
Oleh karena itu,
\begin{equation*}
X(x) =
\frac{F_0}{\omega^2} \left(
\cos \left(  \frac{\omega}{a} x \right) -
\frac{\cos \left( \frac{\omega L}{a} \right) - 1}%
{\sin \left( \frac{\omega L}{a} \right)}
\sin \left( \frac{\omega}{a} x \right)
- 1
\right) .
\end{equation*}
Solusi khusus $y_p$ yang kita cari adalah
\begin{equation*}
\mybxbg{~~
y_p(x,t) =
\frac{F_0}{\omega^2} \left(
\cos \left( \frac{\omega}{a} x \right) -
\frac{\cos \left( \frac{\omega L}{a} \right) - 1}%
{\sin \left( \frac{\omega L}{a}\right)}
\sin \left( \frac{\omega}{a} x \right)
-1
\right)
\cos ( \omega t) .
~~}
\end{equation*}

\begin{exercise}
Periksa bahwa $y_p$ bekerja.
\end{exercise}

Kita sampai pada pokok yang sebelumnya kita lewatkan.  Misalkan 
$\sin ( \frac{\omega L}{a} ) = 0$.  Ini berarti bahwa
$\omega$ sama dengan salah satu frekuensi alami sistem,
yaitu suatu kelipatan frekuensi dasar $\frac{\pi a}{L}$.
Jika $\omega$
tidak sama dengan suatu kelipatan $\frac{\pi a}{L}$,
tetapi sangat dekat,
maka koefisien $B$ dalam \eqref{natfreq:Beq} tampaknya
menjadi sangat besar ketika penyebutnya menuju nol.
Namun, janganlah kita tergesa-gesa menarik kesimpulan
karena pembilangnya juga mungkin menuju nol.
Ketika $\omega = \frac{n \pi a}{L}$
untuk $n$ genap, maka $\cos (\frac{\omega L}{a}) = 1$.
Kita dapat mengambil
$B=0$ untuk memperoleh solusi terbatas yang baik dengan bentuk yang sama,
$y_p = 
\frac{F_0}{\omega^2}
\bigl(
 \cos\bigl(\frac{\omega}{a} x\bigr) -  1 \bigr)
\cos(\omega t)$, sehingga tidak terjadi resonansi.
Resonansi hanya terjadi ketika 
baik $\cos (\frac{\omega L}{a}) = -1$ maupun
$\sin (\frac{\omega L}{a}) = 0$,
yaitu ketika $\omega = \frac{n \pi a }{L}$
untuk $n$ \emph{ganjil}.
Ketika $n$ ganjil, jika kita mengambil $\omega$ yang mendekati
$\frac{n \pi a }{L}$, maka pembilang dalam koefisien
$B$ tidak menuju nol, tetapi penyebutnya menuju nol, sehingga $B$ benar-benar
\myquote{meledak.}
Kita dapat kembali mencoba menyelesaikan secara eksplisit solusi resonansinya jika menginginkannya, tetapi dalam pengertian yang tepat, solusi tersebut adalah limit solusi-solusi ketika $\omega$
mendekati suatu frekuensi resonansi.
Dalam kehidupan nyata, resonansi murni bagaimanapun tidak pernah terjadi.
%Anda akan memperoleh solusi khusus yang memuat suatu suku
%seperti $t \cos (\omega t)$ (meskipun lebih sulit, tetapi masih mungkin,
%untuk mencari solusinya dengan metode di atas).

Perhitungan frekuensi resonansi di atas menjelaskan mengapa suatu dawai mulai
bergetar jika
dawai yang identik dipetik di dekatnya.  Tanpa adanya gesekan, getaran ini
akan semakin keras seiring berjalannya waktu.
Di sisi lain, kecil kemungkinan Anda memperoleh getaran besar jika frekuensi pemaksa 
tidak dekat dengan frekuensi resonansi, bahkan jika terdapat mesin jet yang
berjalan di dekat
dawai tersebut.  Artinya, amplitudo tidak terus
meningkat kecuali Anda menyetelnya tepat pada frekuensi yang sesuai.

Fenomena resonansi serupa terjadi ketika Anda memecahkan gelas anggur dengan suara manusia
(ya,
hal ini mungkin dilakukan, tetapi tidak mudah%
\footnote{\emph{Mythbusters}, episode 31, Discovery Channel, pertama kali ditayangkan
18 Mei 2005.}) jika Anda kebetulan mencapai frekuensi yang tepat.
Namun, gelas memiliki bunyi yang jauh lebih murni, yaitu lebih menyerupai
vibrafon, sehingga jauh lebih sedikit frekuensi resonansi yang dapat dicapai.

Ketika fungsi pemaksa lebih rumit, Anda mendekomposisinya dalam
deret Fourier dan menerapkan hasil di atas.  Anda mungkin juga perlu menyelesaikan
masalah di atas jika fungsi pemaksanya berupa sinus alih-alih kosinus,
tetapi jika dipikirkan, solusinya hampir sama.
%  Artinya,
%solusi khususnya memuat $\sin$ alih-alih kosinus dan selebihnya identik.
%Solusi komplemennya akan sedikit berbeda karena $f$ dan $g$
%akan berbeda.

\begin{example}
%Mari lakukan perhitungannya untuk nilai-nilai tertentu.
Misalkan $F_0 = 1$, $\omega = 1$, $L=1$, dan $a=1$.
Mari kita cari suatu solusi $y$ untuk \eqref{sps:forcedeq}
yang memenuhi syarat batas nol
\eqref{sps:forcedinitcond}.  Pertama, kita memperoleh
\begin{equation*}
y_p(x,t) =
\left(
\cos (x) -
\frac{\cos (1) - 1}{\sin (1)}
\sin (x)
-1
\right)
\cos (t) .
\end{equation*}
Tuliskan $B = \frac{\cos (1) - 1}{\sin (1)}$ demi kesederhanaan.

Substitusikan $t=0$ ke dalam $y_p$ untuk memperoleh
\begin{equation*}
f(x) =- y_p(x,0) = 
- \cos x +
B \sin x
+1 ,
\end{equation*}
dan setelah menurunkan $y_p$ terhadap $t$, kita melihat bahwa 
$g(x) = -\frac{\partial y_p}{\partial t}(x,0) = 0$.

Oleh karena itu, untuk mencari $y_c$, kita perlu menyelesaikan masalah
\begin{align*}
& w_{tt} = w_{xx} , \\
& w(0,t) = 0 , \quad w(1,t) = 0 , \\
& w(x,0) = f(x) = - \cos x + B \sin x +1 , \\
& w_t(x,0) = g(x) = 0 ,
\end{align*}
dan menetapkan $y_c=w$.
Jika kita menggunakan d'Alembert, kita perhatikan bahwa
rumus yang digunakan untuk mendefinisikan $w(x,0)$ tidak ganjil,
sehingga tidak cukup sekadar menyubstitusikan ekspresi $w(x,0)$
langsung ke dalam rumus
d'Alembert!  Anda harus mendefinisikan $F$ sebagai perluasan ganjil yang periodik-2
dari $w(x,0)$.  Kemudian solusi kita $y = y_p+y_c$ untuk
\eqref{sps:forcedeq} yang tunduk pada \eqref{sps:forcedinitcond}
adalah
\begin{equation} \label{natfreq:exsol}
y(x,t) = 
\frac{F(x+t) + F(x-t)}{2} + 
\left(
\cos (x) -
\frac{\cos (1) - 1}{\sin (1)}
\sin (x)
-1
\right)
\cos (t) .
\end{equation}

Tidaklah sulit untuk menghitung nilai-nilai tertentu
bagi perluasan periodik ganjil suatu fungsi, sehingga
\eqref{natfreq:exsol} merupakan solusi yang sangat baik untuk masalah tersebut.
Sebagai contoh, sangat mudah meminta komputer melakukannya, tidak seperti solusi deret.
Suatu grafik diberikan dalam \figurevref{natfreq:forcedvibfig}.
\begin{myfig}
\capstart
\diffyincludegraphics{width=5in}{width=7.5in}{natfreq-forcedvib}{%
Grafik suatu permukaan di atas bidang xt, dengan x berkisar dari 0 hingga 1 dan t
berkisar dari 0 hingga 5.  Fungsi y bernilai nol pada 3 sisi, yaitu pada t sama dengan 0
dan ketika x sama dengan 0 atau x sama dengan 1.  Di bagian dalam domainnya, ketika t membesar, terdapat suatu tonjolan yang naik
hingga sedikit di atas 0,2, lalu terdapat dua tonjolan yang turun ke bawah bidang xt,
dengan tonjolan kedua tampak sedikit lebih menonjol.
Kemudian fungsi mulai membentuk tonjolan lain ke atas.}
\caption{Grafik $y(x,t) = \frac{F(x+t) + F(x-t)}{2} + \left( \cos (x) -
\frac{\cos (1) - 1}{\sin (1)} \sin (x) -1 \right) \cos (t)$.%
\label{natfreq:forcedvibfig}}
\end{myfig}
\end{example}

\subsection{Osilasi suhu bawah tanah}

Misalkan $u(x,t)$ adalah suhu di suatu lokasi pada kedalaman $x$
di bawah tanah pada waktu $t$.  Lihat \figurevref{sps:groundtempfig}.

\begin{mywrapfig}{2.65in}
\capstart
\inputpdft{sps-groundtemp}{%
Diagram permukaan tanah dengan seseorang berdiri di atasnya, dan sebuah panah
ke bawah yang ditandai dengan `kedalaman x' menunjukkan bahwa x berarti kedalaman di bawah
permukaan tanah.}
\caption{Suhu bawah tanah.\label{sps:groundtempfig}}
\end{mywrapfig}

Suhu $u$ memenuhi persamaan panas $u_t = ku_{xx}$, dengan $k$
merupakan difusivitas tanah.
Kita mengetahui suhu di permukaan $u(0,t)$ dari catatan
cuaca.  Demi kesederhanaan, mari kita asumsikan bahwa
\begin{equation*}
u(0,t) = T_0 + A_0 \cos (\omega t) ,
\end{equation*}
dengan $T_0$ merupakan suhu rata-rata
tahunan, dan
$t=0$ merupakan pertengahan musim panas (tambahkan
tanda negatif di atas untuk menjadikannya pertengahan musim dingin jika Anda menginginkannya).
Konstanta $A_0$ memberikan 
variasi khas sepanjang tahun:
Suhu terpanas adalah $T_0 + A_0$ dan suhu terdingin adalah $T_0 - A_0$.
Demi kesederhanaan, kita mengasumsikan bahwa $T_0 = 0$.
Frekuensi $\omega$ dipilih bergantung pada satuan $t$, sedemikian sehingga
ketika $t=\unit[1]{tahun}$, maka $\omega t = 2 \pi$.  Sebagai contoh, jika $t$
dinyatakan dalam tahun, maka $\omega = 2\pi$.

Tampaknya masuk akal bahwa suhu pada kedalaman $x$ juga berosilasi
dengan frekuensi yang sama.  Hal ini sebenarnya merupakan solusi periodik
tunak, yaitu solusi yang tidak bergantung pada syarat-syarat awal.
Jadi, kita mencari solusi berbentuk
\begin{equation*}
u(x,t) = V(x) \cos (\omega t) + W (x) \sin ( \omega t)
\end{equation*}
untuk masalah
\begin{equation} \label{sps:ueq}
u_t = k u_{xx}, \qquad u(0,t) = A_0 \cos ( \omega t) .
\end{equation}

Kita menggunakan eksponensial kompleks di sini untuk menyederhanakan perhitungan.
Misalkan kita memiliki suatu fungsi bernilai kompleks
\begin{equation*}
h(x,t) = X(x)\, e^{i\omega t} .
\end{equation*}
Kita mencari $h$ sedemikian sehingga $\operatorname{Re} h = u$.
Untuk mencari $h$ yang bagian realnya memenuhi \eqref{sps:ueq}, kita mencari
$h$ sedemikian sehingga
\begin{equation} \label{sps:heq}
h_t = k h_{xx}, \qquad h(0,t) = A_0 e^{i\omega t} .
\end{equation}

\begin{exercise}
Misalkan $h$ memenuhi \eqref{sps:heq}.
Gunakan \hyperref[eulersformula]{rumus Euler} untuk eksponensial kompleks guna
memeriksa bahwa $u = \operatorname{Re} h$ memenuhi \eqref{sps:ueq}.
\end{exercise}

Substitusikan $h$ ke dalam \eqref{sps:heq}.
\begin{equation*}
i\omega X e^{i\omega t} = k X'' e^{i \omega t} .
\end{equation*}
Oleh karena itu,
\begin{equation*}
k X''  - i \omega X = 0 ,
\end{equation*}
atau 
\begin{equation*}
X''  - \alpha^2 X = 0 ,
\end{equation*}
dengan $\alpha = \pm \sqrt{\frac{i\omega}{k}}$.  Perhatikan bahwa $\pm \sqrt{i} = \pm
\frac{1+i}{\sqrt{2}}$, sehingga Anda dapat menyederhanakannya menjadi
$\alpha = \pm (1+i)\sqrt{\frac{\omega}{2k}}$.
Oleh karena itu, solusi umumnya adalah
\begin{equation*}
X(x) = A e^{-(1+i)\sqrt{\frac{\omega}{2k}} \, x}
+ B e^{(1+i)\sqrt{\frac{\omega}{2k}} \, x} .
\end{equation*}
Kita mengasumsikan bahwa suatu $X(x)$ yang menyelesaikan masalah tersebut harus terbatas ketika $x \to
\infty$, karena $u(x,t)$ seharusnya terbatas
(kita tidak sedang memikirkan inti Bumi!).
Jika Anda menggunakan \hyperref[eulersformula]{rumus Euler} untuk mengembangkan eksponensial kompleks, 
perhatikan bahwa suku kedua takterbatas (jika $B \neq 0$),
sedangkan suku pertama selalu terbatas.  Oleh karena itu, $B=0$.

\begin{exercise}
Gunakan \hyperref[eulersformula]{rumus Euler} untuk menunjukkan bahwa
$e^{(1+i)\sqrt{\frac{\omega}{2k}} \, x}$ takterbatas ketika $x \to \infty$,
sedangkan $e^{-(1+i)\sqrt{\frac{\omega}{2k}} \, x}$ terbatas
ketika $x \to \infty$.
\end{exercise}

Selain itu, $X(0) = A_0$ karena $h(0,t) = A_0 e^{i \omega t}$.
Dengan demikian, $A=A_0$.  Ini berarti bahwa
\begin{equation*}
h(x,t) = A_0 e^{-(1+i)\sqrt{\frac{\omega}{2k}} \, x} e^{i \omega t}
=
A_0 e^{-(1+i)\sqrt{\frac{\omega}{2k}} \, x + i \omega t}
=
A_0 e^{-\sqrt{\frac{\omega}{2k}} \, x}
e^{i(\omega t - \sqrt{\frac{\omega}{2k}} \, x)} .
\end{equation*}
Kita perlu mengambil bagian real dari $h$, sehingga
kita menerapkan \hyperref[eulersformula]{rumus Euler} untuk memperoleh
\begin{equation*}
h(x,t) =
A_0 e^{-\sqrt{\frac{\omega}{2k}} \, x}
\left(\cos \left(\omega t - \sqrt{\frac{\omega}{2k}}\, x\right) + 
i \sin \left(\omega t - \sqrt{\frac{\omega}{2k}}\, x\right) \right) .
\end{equation*}
Kemudian, akhirnya,
\begin{equation*}
u(x,t) = \operatorname{Re} h(x,t) =
A_0 e^{-\sqrt{\frac{\omega}{2k}}\, x}
\cos \left(\omega t - \sqrt{\frac{\omega}{2k}}\, x\right) .
\end{equation*}
Berhasil!

Perhatikan bahwa fasenya berbeda pada kedalaman yang berbeda.  Pada kedalaman $x$,
fasenya tertunda sebesar $x \sqrt{\frac{\omega}{2k}}$.
Sebagai contoh, dalam satuan cgs (sentimeter-gram-detik)\index{satuan cgs},
kita memiliki $k=0.005$ (nilai khas untuk tanah),
$\omega = \frac{2\pi}{\text{detik dalam setahun}}
= \frac{2\pi}{31,557,341} \approx 1.99 \times {10}^{-7}$.   Kemudian,
jika kita menghitung tempat pergeseran fase $x \sqrt{\frac{\omega}{2k}} = \pi$,
kita menemukan kedalaman dalam sentimeter tempat musim-musim terbalik.  Artinya,
kita memperoleh kedalaman tempat musim panas menjadi yang terdingin dan musim dingin menjadi yang terhangat.
Kita memperoleh
sekitar 700 sentimeter, yaitu sekitar 23 kaki di bawah permukaan tanah.

Berhati-hatilah agar tidak tergesa-gesa menarik kesimpulan.  Ayunan suhu meluruh dengan cepat saat Anda menggali lebih dalam.  Amplitudo
ayunan suhu adalah
$A_0 e^{-\sqrt{\frac{\omega}{2k}} x}$.  Fungsi ini meluruh \emph{sangat}
cepat ketika $x$ (kedalaman) bertambah.
Mari kita kembali mengambil
parameter khas seperti di atas.  Kita juga mengasumsikan bahwa
ayunan suhu permukaan kita adalah $\pm {15}^\circ$ Celsius, yaitu
$A_0 = 15$.  Maka variasi suhu maksimum pada kedalaman 700 sentimeter
hanya $\pm {0.66}^\circ$ Celsius.

Anda tidak perlu menggali sangat dalam untuk memperoleh
\myquote{lemari pendingin} yang efektif, dengan suhu hampir konstan.  Itulah sebabnya anggur disimpan di ruang bawah tanah; Anda memerlukan
suhu yang konsisten.
Perbedaan suhu juga dapat dimanfaatkan sebagai energi.  Sebuah rumah dapat
dipanaskan atau didinginkan dengan memanfaatkan fakta di atas.
Bahkan tanpa inti Bumi, Anda dapat memanaskan rumah pada musim dingin dan mendinginkannya
pada musim panas.  Inti Bumi membuat
suhu semakin tinggi semakin dalam Anda menggali, meskipun Anda perlu menggali cukup
dalam untuk merasakan perbedaannya.
Kita tidak memperhitungkan hal itu di atas.

\subsection{Latihan}

\begin{exercise} \label{sps:sinforceexr}
Misalkan fungsi pemaksa untuk dawai bergetar
adalah $F_0 \sin (\omega t)$.  Turunkan solusi khusus $y_p$.
\end{exercise}

\begin{exercise}
Perhatikan dawai bergetar paksa.
Misalkan $L=1$, $a=1$.  Misalkan fungsi pemaksanya
adalah gelombang persegi yang bernilai 1 pada interval $0 < t < \pi$ dan
$-1$ pada interval $-\pi < t < 0$.
Carilah solusi khususnya.  Petunjuk: Anda mungkin ingin menggunakan hasil
dari \exerciseref{sps:sinforceexr}.
\end{exercise}

\begin{exercise}
Satuannya adalah cgs (sentimeter-gram-detik)\index{satuan cgs}.
Untuk $k=0.005$, $\omega = 1.991 \times {10}^{-7}$, $A_0 = 20$.
Carilah kedalaman tempat variasi suhu menjadi setengah ($\pm 10$
derajat) dari variasinya di permukaan.
\end{exercise}

\begin{exercise}
Turunkan solusi untuk osilasi suhu bawah tanah tanpa mengasumsikan
bahwa $T_0 = 0$.
\end{exercise}

\setcounter{exercise}{100}

\begin{exercise}
Perhatikan dawai bergetar paksa.
Misalkan $L=1$, $a=1$.  Misalkan fungsi pemaksanya
merupakan gelombang segitiga, $\lvert t \rvert -\frac{\pi}{2}$
pada $-\pi < t < \pi$ yang diperluas secara periodik.
Carilah solusi khususnya.
\end{exercise}
\exsol{%
%Deret Fourier fungsi pemaksanya adalah:
%\sum\limits_{\substack{n=1 \\ n \text{ ganjil}}}^\infty
%\frac{-4}{\pi n^2}
%\cos (n t)
$y_p(x,t) =
\sum\limits_{\substack{n=1 \\ n \text{ ganjil}}}^\infty
\frac{-4}{\pi n^4}
\left(
\cos(n x ) -
\frac{\cos ( n ) - 1}%
{\sin( n )}
\sin( n x)
-1
\right)
\cos (n t) .
$
}

\begin{exercise}
Satuannya adalah cgs (sentimeter-gram-detik)\index{satuan cgs}.
Untuk $k=0.01$, $\omega = 1.991 \times {10}^{-7}$, $A_0 = 25$.
Carilah kedalaman tempat musim panas kembali menjadi titik terpanas.
\end{exercise}
\exsol{%
%$x \sqrt{\frac{\omega}{2k}} = 2\pi$
%$x = 2\pi\sqrt{\frac{2k}{\omega}}$
%$x = 2\pi\sqrt{\frac{0.02}{1.991 x 10^-7}}$
Sekitar 1991 sentimeter
}
